% ==============================================================================
% Seccion: Plataforma de Adquisicion y Estandarizacion Experimental (Nandu LSD)
% Archivo: software.tex
% Enfoque: Experimental, Fisiologico y Biomedico (No Tecnico en Software)
% Estilo: Espanol Academico FCEyN / UBA (Laboratorio de Sistemas Dinamicos)
% ==============================================================================

\section{Plataforma de Adquisición y Estandarización Experimental}
\label{sec:plataforma_software}

El registro de electromiografía de superficie facial (sEMG) presenta desafíos metrológicos y fisiológicos severos. Los biopotenciales superficiales generados por la musculatura fonatoria exhiben amplitudes reducidas (típicamente entre $10\,\mu\text{V}$ y $1\,\text{mV}$), alta sensibilidad a artefactos por desplazamiento de contacto piel-electrodo, interferencia electromagnética de red ($50\,\text{Hz}$) y una notable variabilidad intrínseca en la coordinación motora inter e intra-sujeto. Para superar estas fuentes de dispersión y garantizar la reproducibilidad y validez estadística de los ensayos, se desarrolló la plataforma científica \textit{Ñandú LSD}, concebida específicamente para sistematizar los protocolos de captura, curación y organización de registros mioeléctricos en el marco de interfaces de habla silenciosa (\textit{Silent Speech Interfaces}, SSI).

\subsection{Protocolo Automatizado de Guiado y Captura (AutoForge)}
\label{subsec:protocolo_autoforge}

Uno de los factores determinantes en la dispersión de las mediciones sEMG es la falta de sincronismo temporal y la fatiga neuromuscular inducida por repeticiones no pautadas. El sistema incorpora un protocolo de adquisición automatizado (\textit{AutoForge}) que actúa como un director experimental estructurado:

\begin{enumerate}
    \item \textbf{Calibración Dinámica de Ruido Basal:} Antes de cada ciclo de fonación o gesto motor, el sistema muestrea de forma autónoma una ventana temporal de reposo eléctrico previo al estímulo. Esto permite caracterizar la línea base del sujeto en tiempo real y calcular umbrales estadísticos de activación muscular ($\mu_{\text{ruido}} + k \cdot \sigma_{\text{ruido}}$) adaptados a la impedancia de contacto electrodo-tejido en ese instante específico.
    \item \textbf{Pauta Temporal por Metrónomo Audiovisual:} La contracción y la relajación del sujeto son guiadas mediante estímulos visuales de alto contraste y marcas acústicas periódicas a una frecuencia programable (típicamente entre $20$ y $40\,\text{BPM}$). Este diseño fija una ventana de expectativa temporal constante, reduce drásticamente el \textit{jitter} de tiempo de reacción y estandariza la duración y cinemática del gesto articulatorio a lo largo de decenas de repeticiones.
    \item \textbf{Grabación de Secuencias Continuas y Ciclado de Fonemas:} El protocolo permite iterar de manera automatizada sobre diccionarios predefinidos de vocales y palabras (e.g., /a/, /e/, /i/, /o/, /u/), asociando a cada ventana de activación su correspondiente etiqueta fonética y metadatos contextuales. Esto elimina por completo la necesidad de intervención manual del operador durante el ensayo, evitando sesgos de manipulación y distracciones en el participante.
\end{enumerate}

\subsection{Preservación de la Sinergia Biomecánica y Alineación Master-Slave}
\label{subsec:sinergia_master_slave}

La producción de habla, tanto sonora como silenciosa, no se origina en la acción aislada de un único efector muscular, sino en la co-activación coordinada y espaciotemporal de complejos musculares faciales y suprahioideos (\textit{Orbicularis Oris}, \textit{Depressor Anguli Oris}, \textit{Mylohyoid}). En este marco biomecánico, la diferencia temporal relativa (\textit{phase lag}) entre el encendido de los distintos grupos musculares constituye la firma fisiológica determinante para discriminar fonemas que presentan niveles de energía muscular globalmente similares.

Para evitar la pérdida de esta correlación de fase temporal durante el recorte y preprocesamiento, la plataforma implementa una estrategia de sincronización \textit{Master-Slave}:
\begin{itemize}
    \item Se establece un canal maestro de referencia temporal (pudiendo corresponder a la envolvente de la señal acústica de control o al canal muscular primario con mayor relación señal-ruido) que determina de forma no supervisada el instante de ignición o de máxima activación del gesto fonatorio ($t_0$).
    \item Todos los canales musculares esclavos se recortan y sincronizan tomando rígidamente la misma ventana temporal $[t_0 - \Delta t_{\text{pre}}, t_0 + \Delta t_{\text{post}}]$ dictada por el canal maestro.
\end{itemize}
De esta forma, se conservan íntegramente los desfases fisiológicos relativos inter-musculares, impidiendo la introducción de deformaciones artificiales en las matrices de correlación cruzada, en las proyecciones en componentes principales (PCA) y en las variedades topológicas no lineales (UMAP).

\subsection{Control de Calidad en Línea y Monitoreo de Fatiga Muscular In-Situ}
\label{subsec:control_calidad}

La adquisición prolongada de bioseñales mioeléctricas se encuentra expuesta a la degradación progresiva de la calidad de registro por factores como sudoración, micromovimientos del parche de electrodos o agotamiento neuromuscular del sujeto. La plataforma integra herramientas de diagnóstico en tiempo real para el control de calidad:
\begin{itemize}
    \item \textbf{Evaluación Instantánea de la Relación Señal-Ruido (SNR):} Para cada pulso mioeléctrico detectado, se calcula el cociente energético entre la ventana de contracción activa y la ventana basal de reposo previa:
    \begin{equation}
        \text{SNR}_{\text{pulso}} = 10 \log_{10} \left( \frac{E_{\text{contracción}}}{E_{\text{reposo}}} \right) = 10 \log_{10} \left( \frac{\sum_{n \in \mathcal{W}_{\text{act}}} s^2[n]}{\frac{|\mathcal{W}_{\text{act}}|}{|\mathcal{W}_{\text{rep}}|} \sum_{m \in \mathcal{W}_{\text{rep}}} s^2[m]} \right)
        \label{eq:snr_pulso}
    \end{equation}
    Si este valor desciende por debajo del umbral de calidad requerido, el ensayo se etiqueta para revisión o descarte automático.
    \item \textbf{Monitoreo del Ruido Inter-pulso (Evaluador de Relajación):} El sistema cuantifica la densidad de energía remanente en el punto medio del intervalo de descanso entre contracciones consecutivas. Un incremento sostenido del valor eficaz ($\text{RMS}$) en el reposo inter-pulso actúa como biomarcador de fatiga neuromuscular o de incapacidad de retorno a la línea base, permitiendo al operador pausar la sesión antes de incorporar datos contaminados por espasmos o tono basal elevado.
\end{itemize}

\subsection{Organización Jerárquica de Datos y Reproducibilidad Experimental}
\label{subsec:organizacion_datos}

Para garantizar la rigurosa reproducibilidad del protocolo experimental y alimentar de forma determinística los modelos de aprendizaje automático y decodificación neuronal, los registros se estructuran jerárquicamente en el almacenamiento:
\begin{itemize}
    \item Cada sesión de laboratorio se encapsula en una estructura unívoca indexada por fecha de adquisición y denominación del protocolo experimental (\nolinkurl{base_de_datos_electrodos/<Fecha>/<Sesion>/}).
    \item Cada canal físico ($0, 1, 2, 3$) posee un subdirectorio aislado (\texttt{canal\_0}, \texttt{canal\_1}, \texttt{canal\_2}, \texttt{canal\_3}) que preserva la señal mioeléctrica continua sin pérdidas y un archivo descriptivo de metadatos en formato JSON.
    \item Los metadatos encapsulan las condiciones experimentales completas: frecuencia de digitalización ($f_s$), ganancia analógica de instrumentación ($\sim 500\times$), resistencia estimada de electrodos ($R_{\text{ohm}}$), cadencia del metrónomo (BPM), diccionario de fonemas articulados y asignación anatómica de músculos.
\end{itemize}
Esta disposición asegura que las etapas posteriores de análisis (filtrado digital de fase cero, extracción de envolventes RMS, segmentación y conformación de tensores para redes neuronales) sean reproducibles y trazables a lo largo del tiempo.

En el Cuadro~\ref{tab:dimensiones_plataforma} se sintetizan los pilares metodológicos de la plataforma y su impacto directo en el procesamiento y modelado biomédico.

\begin{table}[htbp]
\centering
\small
\caption{Dimensiones metodológicas de la plataforma \textit{Ñandú LSD} y su impacto en la adquisición de sEMG para interfaces de habla silenciosa (SSI).}
\label{tab:dimensiones_plataforma}
\vspace{0.2cm}
\begin{tabular}{p{2.8cm} p{5.2cm} p{6.0cm}}
\toprule
\textbf{Eje Metodológico} & \textbf{Estrategia Experimental} & \textbf{Impacto en el Modelado de SSI} \\
\midrule
Sincronismo temporal de activación & Metrónomo audiovisual pautado ($20\text{ a }40\,\text{BPM}$) en protocolo AutoForge & Reducción drástica del \textit{jitter} de reacción y homogeneidad cinemática del gesto articulado. \\
\addlinespace
Calibración dinámica de línea base & Muestreo previo de reposo eléctrico y umbral adaptativo ($\mu + k\sigma$) & Estimación de $\text{SNR}$ instantánea y robustez frente a derivas de impedancia de contacto. \\
\addlinespace
Preservación de sinergia muscular & Segmentación \textit{Master-Slave} anclada al canal rector ($t_0$) & Conservación de retardos relativos de fase fisiológicos (\textit{lags}) para PCA, UMAP y clasificadores. \\
\addlinespace
Monitoreo de fatiga \textit{in-situ} & Evaluación del valor eficaz ($\text{RMS}$) en el reposo inter-pulso & Detección preventiva de agotamiento motor y descarte objetivo de registros contaminados. \\
\addlinespace
Trazabilidad y reproducibilidad & Jerarquía de carpetas multicanal estandarizada y metadatos JSON unificados & Ingesta determinística y estructuración de tensores para autoencoders 1D y modelos neuronales. \\
\bottomrule
\end{tabular}
\end{table}

\subsection{Relevancia Biomédica e Impacto Clínico en Habla Silenciosa}
\label{subsec:relevancia_biomedica}

El diseño de la plataforma responde a un objetivo de investigación biomédica de alto impacto: proveer un marco experimental robusto para el desarrollo de sistemas de comunicación aumentativa y alternativa destinados a personas con afonía severa, parálisis de cuerdas vocales, pacientes laringectomizados o con enfermedades de la motoneurona (como la Esclerosis Lateral Amiotrófica, ELA). Al capturar fielmente las dinámicas electromiográficas en ausencia de emisión sonora audible (\textit{silent speech}), el sistema proporciona las bases de datos de alta fidelidad necesarias para mapear la intención articulatoria periférica hacia modelos generativos de habla sintética.

